% Options for packages loaded elsewhere
\PassOptionsToPackage{unicode}{hyperref}
\PassOptionsToPackage{hyphens}{url}
\PassOptionsToPackage{dvipsnames,svgnames,x11names}{xcolor}
%
\documentclass[
  11pt,
]{article}

\usepackage{amsmath,amssymb}
\usepackage{iftex}
\ifPDFTeX
  \usepackage[T1]{fontenc}
  \usepackage[utf8]{inputenc}
  \usepackage{textcomp} % provide euro and other symbols
\else % if luatex or xetex
  \usepackage{unicode-math}
  \defaultfontfeatures{Scale=MatchLowercase}
  \defaultfontfeatures[\rmfamily]{Ligatures=TeX,Scale=1}
\fi
\usepackage{lmodern}
\ifPDFTeX\else  
    % xetex/luatex font selection
    \setmainfont[]{Arial}
\fi
% Use upquote if available, for straight quotes in verbatim environments
\IfFileExists{upquote.sty}{\usepackage{upquote}}{}
\IfFileExists{microtype.sty}{% use microtype if available
  \usepackage[]{microtype}
  \UseMicrotypeSet[protrusion]{basicmath} % disable protrusion for tt fonts
}{}
\makeatletter
\@ifundefined{KOMAClassName}{% if non-KOMA class
  \IfFileExists{parskip.sty}{%
    \usepackage{parskip}
  }{% else
    \setlength{\parindent}{0pt}
    \setlength{\parskip}{6pt plus 2pt minus 1pt}}
}{% if KOMA class
  \KOMAoptions{parskip=half}}
\makeatother
\usepackage{xcolor}
\usepackage[a4paper, margin=0.8in]{geometry}
\setlength{\emergencystretch}{3em} % prevent overfull lines
\setcounter{secnumdepth}{5}
% Make \paragraph and \subparagraph free-standing
\makeatletter
\ifx\paragraph\undefined\else
  \let\oldparagraph\paragraph
  \renewcommand{\paragraph}{
    \@ifstar
      \xxxParagraphStar
      \xxxParagraphNoStar
  }
  \newcommand{\xxxParagraphStar}[1]{\oldparagraph*{#1}\mbox{}}
  \newcommand{\xxxParagraphNoStar}[1]{\oldparagraph{#1}\mbox{}}
\fi
\ifx\subparagraph\undefined\else
  \let\oldsubparagraph\subparagraph
  \renewcommand{\subparagraph}{
    \@ifstar
      \xxxSubParagraphStar
      \xxxSubParagraphNoStar
  }
  \newcommand{\xxxSubParagraphStar}[1]{\oldsubparagraph*{#1}\mbox{}}
  \newcommand{\xxxSubParagraphNoStar}[1]{\oldsubparagraph{#1}\mbox{}}
\fi
\makeatother


\providecommand{\tightlist}{%
  \setlength{\itemsep}{0pt}\setlength{\parskip}{0pt}}\usepackage{longtable,booktabs,array}
\usepackage{calc} % for calculating minipage widths
% Correct order of tables after \paragraph or \subparagraph
\usepackage{etoolbox}
\makeatletter
\patchcmd\longtable{\par}{\if@noskipsec\mbox{}\fi\par}{}{}
\makeatother
% Allow footnotes in longtable head/foot
\IfFileExists{footnotehyper.sty}{\usepackage{footnotehyper}}{\usepackage{footnote}}
\makesavenoteenv{longtable}
\usepackage{graphicx}
\makeatletter
\newsavebox\pandoc@box
\newcommand*\pandocbounded[1]{% scales image to fit in text height/width
  \sbox\pandoc@box{#1}%
  \Gscale@div\@tempa{\textheight}{\dimexpr\ht\pandoc@box+\dp\pandoc@box\relax}%
  \Gscale@div\@tempb{\linewidth}{\wd\pandoc@box}%
  \ifdim\@tempb\p@<\@tempa\p@\let\@tempa\@tempb\fi% select the smaller of both
  \ifdim\@tempa\p@<\p@\scalebox{\@tempa}{\usebox\pandoc@box}%
  \else\usebox{\pandoc@box}%
  \fi%
}
% Set default figure placement to htbp
\def\fps@figure{htbp}
\makeatother

\makeatletter
\@ifpackageloaded{caption}{}{\usepackage{caption}}
\AtBeginDocument{%
\ifdefined\contentsname
  \renewcommand*\contentsname{Table of contents}
\else
  \newcommand\contentsname{Table of contents}
\fi
\ifdefined\listfigurename
  \renewcommand*\listfigurename{List of Figures}
\else
  \newcommand\listfigurename{List of Figures}
\fi
\ifdefined\listtablename
  \renewcommand*\listtablename{List of Tables}
\else
  \newcommand\listtablename{List of Tables}
\fi
\ifdefined\figurename
  \renewcommand*\figurename{Figure}
\else
  \newcommand\figurename{Figure}
\fi
\ifdefined\tablename
  \renewcommand*\tablename{Table}
\else
  \newcommand\tablename{Table}
\fi
}
\@ifpackageloaded{float}{}{\usepackage{float}}
\floatstyle{ruled}
\@ifundefined{c@chapter}{\newfloat{codelisting}{h}{lop}}{\newfloat{codelisting}{h}{lop}[chapter]}
\floatname{codelisting}{Listing}
\newcommand*\listoflistings{\listof{codelisting}{List of Listings}}
\makeatother
\makeatletter
\makeatother
\makeatletter
\@ifpackageloaded{caption}{}{\usepackage{caption}}
\@ifpackageloaded{subcaption}{}{\usepackage{subcaption}}
\makeatother

\usepackage{bookmark}

\IfFileExists{xurl.sty}{\usepackage{xurl}}{} % add URL line breaks if available
\urlstyle{same} % disable monospaced font for URLs
\hypersetup{
  pdftitle={01\_Accession\_phenotypes},
  pdfauthor={Henry Janse van Rensburg},
  colorlinks=true,
  linkcolor={blue},
  filecolor={Maroon},
  citecolor={Blue},
  urlcolor={Blue},
  pdfcreator={LaTeX via pandoc}}


\title{01\_Accession\_phenotypes}
\author{Henry Janse van Rensburg}
\date{2025-06-27}

\begin{document}
\maketitle

\renewcommand*\contentsname{Table of contents}
{
\hypersetup{linkcolor=}
\setcounter{tocdepth}{2}
\tableofcontents
}

\section{Background}\label{background}

We conducted this experiment to study the phenotypic diversity within
\emph{Arabidopsis} accession to BX-growth feedbacks. We screened 410
accessions for their growth feedback on two different soils, BXplus and
BXminus. The phenotypes were measured twice a day for 20 days after
germination.

Materials and methodsWe grew 410 Arabidopsis accessions in triplicate on
BXplus and BXminus soils and phenotyped them twice a day. We extracted
the rosette area using ARADEEPOPSIS. The list of accessions used in this
study is available in the file \texttt{bx\_experiment\_genotypes.csv} in
the input folder.

\textbf{Purpose}: We performed this study to eventually perform GWAS on
the phenotypes.

\textbf{Important results:} We found that the phenotypes were
significantly different between the two soil types, and varied based on
plant genotype.

\textbf{Conclusions:} We concluded that the phenotypic diversity within
Arabidopsis accessions is significant and can be used for further
genetic studies.

\section{Setup}\label{setup}

\begin{verbatim}
\small
\end{verbatim}

\subsection{Load libraries}\label{load-libraries}

\subsection{Import phenotypic data}\label{import-phenotypic-data}

\textbf{Purpose}: Imports the phenotype data as
\texttt{plant\_region\_areas}and the SNPs identifid in the GWAs as
\texttt{tot\_plant\_specific\_snps}from CSV files into R as data frames.

\begin{longtable}[]{@{}rrrrlrr@{}}
\caption{SNP table.}\tabularnewline
\toprule\noalign{}
chrom & pos & log10\_p & pv & Significant & mac & dag \\
\midrule\noalign{}
\endfirsthead
\toprule\noalign{}
chrom & pos & log10\_p & pv & Significant & mac & dag \\
\midrule\noalign{}
\endhead
\bottomrule\noalign{}
\endlastfoot
3 & 21437549 & 8.990193 & 0 & Bonferroni & 5 & 13 \\
3 & 21437549 & 8.952353 & 0 & Bonferroni & 5 & 12 \\
3 & 21437549 & 8.870501 & 0 & Bonferroni & 5 & 14 \\
\end{longtable}

\section{Phenotypes}\label{phenotypes}

\subsection{Col-0 phenotypes}\label{col-0-phenotypes}

\pandocbounded{\includegraphics[keepaspectratio]{01_Accession_phenotypes_files/figure-pdf/Col-0 growth over time-1.pdf}}

\pandocbounded{\includegraphics[keepaspectratio]{01_Accession_phenotypes_files/figure-pdf/Col-0 boxplot-1.pdf}}

\begin{verbatim}
# A tibble: 6 x 7
    age estimate1 estimate2 statistic p.value conf.low conf.high
  <dbl>     <dbl>     <dbl>     <dbl>   <dbl>    <dbl>     <dbl>
1    16    97479.   106100.     -1.31  0.203   -22213.     4971.
2    20   119014.   128217.     -1.08  0.293   -26898.     8492.
3    24   211905.   244196.     -2.12  0.0438  -63612.     -971.
4    27   292747.   349615.     -2.44  0.0224 -104980.    -8757.
5    32   687349    839486.     -2.74  0.011  -266327.   -37948.
6    35   821471.  1009477.     -2.74  0.0109 -328973.   -47040.
\end{verbatim}

\begin{verbatim}
# A tibble: 6 x 8
# Groups:   age [6]
    age term     df         sumsq        meansq statistic p.value tukey_hsd
  <dbl> <chr> <dbl>         <dbl>         <dbl>     <dbl>   <dbl> <list>   
1    16 BX        1    520277350.    520277350.      1.70  0.203  <tibble> 
2    20 BX        1    592857260.    592857260.      1.16  0.292  <tibble> 
3    24 BX        1   7299122223.   7299122223.      4.51  0.0434 <tibble> 
4    27 BX        1  22638184046.  22638184046.      5.94  0.0219 <tibble> 
5    32 BX        1 162020580206. 162020580206.      7.50  0.0110 <tibble> 
6    35 BX        1 247424920289. 247424920289.      7.52  0.0109 <tibble> 
\end{verbatim}

\begin{verbatim}
# A tibble: 6 x 10
# Groups:   age [6]
    age data     anova_model term  contrast   null.value estimate conf.low
  <dbl> <list>   <list>      <chr> <chr>           <dbl>    <dbl>    <dbl>
1    16 <tibble> <aov>       BX    plus-minus          0    8621.   -4951.
2    20 <tibble> <aov>       BX    plus-minus          0    9203.   -8389.
3    24 <tibble> <aov>       BX    plus-minus          0   32291.    1038.
4    27 <tibble> <aov>       BX    plus-minus          0   56868.    8912.
5    32 <tibble> <aov>       BX    plus-minus          0  152137.   37948.
6    35 <tibble> <aov>       BX    plus-minus          0  188006.   47066.
# i 2 more variables: conf.high <dbl>, adj.p.value <dbl>
\end{verbatim}

\begin{verbatim}
Linear mixed model fit by REML. t-tests use Satterthwaite's method [
lmerModLmerTest]
Formula: plant_region_area ~ age * BX + (1 | plant_id)
   Data: Col0_plant_region_areas

REML criterion at convergence: 4378.3

Scaled residuals: 
    Min      1Q  Median      3Q     Max 
-1.5858 -0.7632 -0.1574  0.8222  2.4756 

Random effects:
 Groups   Name        Variance  Std.Dev.
 plant_id (Intercept) 2.874e+09  53608  
 Residual             1.771e+10 133085  
Number of obs: 168, groups:  plant_id, 28

Fixed effects:
             Estimate Std. Error        df t value Pr(>|t|)    
(Intercept) -669706.2    60455.3     161.6 -11.078  < 2e-16 ***
age           40572.7     2217.3     138.0  18.299  < 2e-16 ***
BXplus      -185282.0    85496.7     161.6  -2.167  0.03169 *  
age:BXplus    10122.2     3135.7     138.0   3.228  0.00156 ** 
---
Signif. codes:  0 '***' 0.001 '**' 0.01 '*' 0.05 '.' 0.1 ' ' 1

Correlation of Fixed Effects:
           (Intr) age    BXplus
age        -0.941              
BXplus     -0.707  0.666       
age:BXplus  0.666 -0.707 -0.941
\end{verbatim}

\pandocbounded{\includegraphics[keepaspectratio]{01_Accession_phenotypes_files/figure-pdf/Col-0 boxplot time-1.pdf}}

\subsection{Extract Phenotypes for SNP Rank
7}\label{extract-phenotypes-for-snp-rank-7}

\textbf{Purpose}: Extracts and combines phenotype data based on SNP rank
7 and prepares it as the \texttt{plant\_region\_areas\_SNP} dataframe in
for further analysis.

\begin{longtable}[]{@{}rlrrl@{}}
\caption{Phenotype table for SNP rank 7.}\tabularnewline
\toprule\noalign{}
ACC\_ID & bx & plant\_region\_area & dag & SNP \\
\midrule\noalign{}
\endfirsthead
\toprule\noalign{}
ACC\_ID & bx & plant\_region\_area & dag & SNP \\
\midrule\noalign{}
\endhead
\bottomrule\noalign{}
\endlastfoot
9830 & + & 636 & 1 & FALSE \\
6092 & + & 413 & 1 & FALSE \\
6022 & + & 380 & 1 & FALSE \\
\end{longtable}

\subsection{Extract phenotypes for DAG
20}\label{extract-phenotypes-for-dag-20}

\textbf{Purpose}: Create a new dataframe,
\texttt{plant\_region\_areas\_DAG20}, filters the data to include only
the first 20 days (DAG ≤ 20) for further analysis, calculates the mean
plant region area for each combination of factors, removes duplicates
and irrelevant columns, pivots the data to spread `bx' treatment values
into separate columns, calculates the area difference, and ranks the
genotypes by the area difference.

\begin{longtable}[]{@{}lrrrrrr@{}}
\caption{Phenotype table for DAG 20.}\tabularnewline
\toprule\noalign{}
genotype & genotype\_id & dag & - & + & area\_diff & rank \\
\midrule\noalign{}
\endfirsthead
\toprule\noalign{}
genotype & genotype\_id & dag & - & + & area\_diff & rank \\
\midrule\noalign{}
\endhead
\bottomrule\noalign{}
\endlastfoot
DraIV 3-7 & 5921 & 21 & 29516.83 & 115727.00 & 86210.17 & 1 \\
DraIV 3-7 & 5921 & 20 & 23975.00 & 98518.33 & 74543.33 & 2 \\
Hov1-10 & 6035 & 21 & 15445.25 & 86467.33 & 71022.08 & 3 \\
\end{longtable}

\subsection{Assign BX feedback and Prepare Data for
Heatmap}\label{assign-bx-feedback-and-prepare-data-for-heatmap}

\textbf{Purpose}: Assigns a preference (``plus'' or ``minus'') to each
genotype based on the area difference at day 10 and merges the
preference data back into the main dataframe. We chose DAG10, as this is
where SNPMMF1 showed up in the GWAS analysis.

\subsection{Prepare Data for Heatmap for All
Accessions}\label{prepare-data-for-heatmap-for-all-accessions}

\textbf{Purpose}: Prepares the genotype and time order for the heatmap.

\subsection{Plot Growth Over Time for All
Accessions}\label{plot-growth-over-time-for-all-accessions}

\textbf{Purpose}: Creates a heatmap showing the area differences over
time for all accessions, grouped by preference.

\begin{figure}[H]

{\centering \pandocbounded{\includegraphics[keepaspectratio]{01_Accession_phenotypes_files/figure-pdf/Fig-heatmap-ACCs-1.pdf}}

}

\caption{Heatmap showing the evolution of growth feedback over time for
all accessions, grouped by preference.}

\end{figure}%

\subsection{Top 10 Positive and Negative
Accessions}\label{top-10-positive-and-negative-accessions}

\textbf{Purpose}: Filters the data from (\textbf{?@fig-heatmap-ACCs})
include only the top 20 positive and negative accessions based on area
difference and applies a pseudo-logarithmic transformation to the area
difference for better visualization at early timepoints.

\subsection{Hierarchical Clustering and Heatmap
Preparation}\label{hierarchical-clustering-and-heatmap-preparation}

\textbf{Purpos}e: Transforms the data into a wide format suitable for
clustering and performs hierarchical clustering on the distance matrix.
The reordered groups are then used to arrange the data frame for
visualization. Creates and displays a heatmap showing the growth
feedback over time for the top 20 extreme accessions.

\subsection{Hierarchical Clustering: Dynamic Time Warping
(DTW)}\label{hierarchical-clustering-dynamic-time-warping-dtw}

\subsubsection{Prepare Data}\label{prepare-data}

\subsubsection{Perform Hierarchical
Clustering}\label{perform-hierarchical-clustering}

\subsubsection{Heatmap}\label{heatmap}

\begin{figure}[H]

{\centering \pandocbounded{\includegraphics[keepaspectratio]{01_Accession_phenotypes_files/figure-pdf/FigS1b-heatmap-ACCs20-1.pdf}}

}

\caption{Heatmap showing the evolution of growth feedbacks as Pseudo-log
transformed values over time for the top 20 accessions, grouped by
preference. Hierarchical clustering was performed using the Dynamic Time
Warping method to account for changes over time.}

\end{figure}%

\begin{verbatim}
pdf 
  2 
\end{verbatim}

\subsection{Accession feedback
distribution}\label{accession-feedback-distribution}

\subsubsection{Prepare data}\label{prepare-data-1}

We filter the data for DAG 15, where SNPMMF1 was still showing up.

\textbf{Purpose}: Filters the data for day 15, assigns a preference
(``plus'' or ``minus'') based on area difference, and ranks the
accessions accordingly. We used dag 15 as this was the last day where
SNPMMF1 showed up in the GWAS analysis and where all the accessions were
still prsent.

Plot relative growth feedbacks

\textbf{Purpose}: Creates a bar plot showing the relative growth
feedback for each accession at day 15, ranked and colored by preference.

\begin{figure}[H]

{\centering \pandocbounded{\includegraphics[keepaspectratio]{01_Accession_phenotypes_files/figure-pdf/Fig1c-ACC-Distribution-1.pdf}}

}

\caption{Distribution of growth feedbackas for all 410 accesions used in
the GWAS. The most extreme accessions and Col-0 is indicated on the
plot.}

\end{figure}%

\subsection{Growth of extreme
accessions}\label{growth-of-extreme-accessions}

\subsubsection{IP-Hue-3}\label{ip-hue-3}

Filter for IP-Hue-3 and Plot the Growth Over Time

\textbf{Purpose}: Filters the data to include only the genotype with
genotype\_id 9851 (IP-Hue-3), calculates the mean plant region area for
each combination of bx and dag, removes duplicates and irrelevant
columns, pivots the data to spread `bx' treatment values into separate
columns, calculates the area difference, and ranks the genotypes by the
area difference.

Plot the growth for IP-Hue-3 over time

Plot the growth feedback for IP-Hue-3 over time

\textbf{Purpose}: Calculates the difference in mean plant region area
between the two treatments for each day and creates a bar plot showing
the \emph{daily growth feedback} for IP-Hue-3.

\pandocbounded{\includegraphics[keepaspectratio]{01_Accession_phenotypes_files/figure-pdf/IP-Hue-3-1.pdf}}

\subsubsection{Hadd-3}\label{hadd-3}

Filter for Hadd-3 and Plot the Growth Over Time

\textbf{Purpose}: Filters the data to include only the genotype with
genotype\_id 9392 (Hadd-3), calculates the mean plant region area for
each combination of bx and dag, removes duplicates and irrelevant
columns, pivots the data to spread `bx' treatment values into separate
columns, calculates the area difference, and ranks the genotypes by the
area difference.

Plot the growth for Hadd-3 over time

Plot the growth feedback for Hadd-3 over time

\textbf{Purpose}: Calculates the difference in mean plant region area
between the two treatments for each day and creates a bar plot showing
the \emph{daily growth feedback} for Hadd-3.

\pandocbounded{\includegraphics[keepaspectratio]{01_Accession_phenotypes_files/figure-pdf/Hadd-3-1.pdf}}

\section{Geographic distribution}\label{geographic-distribution}

\subsection{Import data}\label{import-data}

\begin{longtable}[]{@{}rllrr@{}}
\caption{Coordinate table.}\tabularnewline
\toprule\noalign{}
ACC\_ID & Country & Region & Lat & Long \\
\midrule\noalign{}
\endfirsthead
\toprule\noalign{}
ACC\_ID & Country & Region & Lat & Long \\
\midrule\noalign{}
\endhead
\bottomrule\noalign{}
\endlastfoot
5921 & CZE & DraIV 3 & 49.4112 & 16.2815 \\
5867 & SWE & Dra2 & 55.7600 & 14.1200 \\
10006 & UKR & Kastel Mountain & 44.6419 & 34.3814 \\
\end{longtable}

\subsection{Correlation of growth feedback and geographical
location}\label{correlation-of-growth-feedback-and-geographical-location}

\subsubsection{Prepare data for correlation and regression
analysis}\label{prepare-data-for-correlation-and-regression-analysis}

\textbf{Purpose:} Filters the data to include only observations from day
10, assigns a preference based on the area difference, sorts by area
difference, and ranks the accessions.

Merge the data tables

\textbf{Purpose:} Merges the phenotype data
(plant\_region\_areas\_DAG10) with the coordinates data
(ACC\_coordinates) based on genotype IDs.

\subsubsection{\texorpdfstring{\textbf{Pearson correlation} and
Regression
Analysis}{Pearson correlation and Regression Analysis}}\label{pearson-correlation-and-regression-analysis}

\begin{longtable}[]{@{}llrrrr@{}}
\caption{Correlation Test Results for Latitude and
Longitude}\tabularnewline
\toprule\noalign{}
& Variable & Correlation & p\_value & CI\_lower & CI\_upper \\
\midrule\noalign{}
\endfirsthead
\toprule\noalign{}
& Variable & Correlation & p\_value & CI\_lower & CI\_upper \\
\midrule\noalign{}
\endhead
\bottomrule\noalign{}
\endlastfoot
cor\ldots1 & Latitude & 0.002 & 0.9751 & -0.096 & 0.099 \\
cor\ldots2 & Longitude & 0.058 & 0.2398 & -0.039 & 0.155 \\
\end{longtable}

\subsubsection{\texorpdfstring{\textbf{Linear regression
model}}{Linear regression model}}\label{linear-regression-model}

\begin{verbatim}

Call:
lm(formula = area_diff ~ Lat + genotype_id, data = coordinates_merged_data)

Residuals:
   Min     1Q Median     3Q    Max 
-29305  -4220    489   4132  47752 

Coefficients:
              Estimate Std. Error t value Pr(>|t|)
(Intercept) -948.88265 3679.18872  -0.258    0.797
Lat           -0.86992   62.11497  -0.014    0.989
genotype_id   -0.03179    0.18395  -0.173    0.863

Residual standard error: 8755 on 405 degrees of freedom
Multiple R-squared:  7.612e-05, Adjusted R-squared:  -0.004862 
F-statistic: 0.01542 on 2 and 405 DF,  p-value: 0.9847
\end{verbatim}

\begin{verbatim}
[1] "Latitude does not significantly affect area_diff."
\end{verbatim}

\begin{verbatim}

Call:
lm(formula = area_diff ~ Long + genotype_id, data = coordinates_merged_data)

Residuals:
   Min     1Q Median     3Q    Max 
-29382  -4218    614   4153  48129 

Coefficients:
             Estimate Std. Error t value Pr(>|t|)
(Intercept) -499.8901  1509.2140  -0.331    0.741
Long          18.4793    14.2306   1.299    0.195
genotype_id   -0.1080     0.1871  -0.577    0.564

Residual standard error: 8737 on 405 degrees of freedom
Multiple R-squared:  0.004222,  Adjusted R-squared:  -0.0006958 
F-statistic: 0.8585 on 2 and 405 DF,  p-value: 0.4246
\end{verbatim}

\begin{verbatim}
[1] "Longitue does not significantly affect area_diff."
\end{verbatim}

\begin{longtable}[]{@{}
  >{\raggedright\arraybackslash}p{(\linewidth - 12\tabcolsep) * \real{0.0986}}
  >{\raggedright\arraybackslash}p{(\linewidth - 12\tabcolsep) * \real{0.1408}}
  >{\raggedleft\arraybackslash}p{(\linewidth - 12\tabcolsep) * \real{0.1408}}
  >{\raggedleft\arraybackslash}p{(\linewidth - 12\tabcolsep) * \real{0.1127}}
  >{\raggedleft\arraybackslash}p{(\linewidth - 12\tabcolsep) * \real{0.1408}}
  >{\raggedleft\arraybackslash}p{(\linewidth - 12\tabcolsep) * \real{0.1972}}
  >{\raggedleft\arraybackslash}p{(\linewidth - 12\tabcolsep) * \real{0.1690}}@{}}
\caption{Linear Regression Results for Latitude and
Longitude}\tabularnewline
\toprule\noalign{}
\begin{minipage}[b]{\linewidth}\raggedright
\end{minipage} & \begin{minipage}[b]{\linewidth}\raggedright
Variable
\end{minipage} & \begin{minipage}[b]{\linewidth}\raggedleft
Intercept
\end{minipage} & \begin{minipage}[b]{\linewidth}\raggedleft
Slope
\end{minipage} & \begin{minipage}[b]{\linewidth}\raggedleft
R\_squared
\end{minipage} & \begin{minipage}[b]{\linewidth}\raggedleft
p\_value\_slope
\end{minipage} & \begin{minipage}[b]{\linewidth}\raggedleft
F\_statistic
\end{minipage} \\
\midrule\noalign{}
\endfirsthead
\toprule\noalign{}
\begin{minipage}[b]{\linewidth}\raggedright
\end{minipage} & \begin{minipage}[b]{\linewidth}\raggedright
Variable
\end{minipage} & \begin{minipage}[b]{\linewidth}\raggedleft
Intercept
\end{minipage} & \begin{minipage}[b]{\linewidth}\raggedleft
Slope
\end{minipage} & \begin{minipage}[b]{\linewidth}\raggedleft
R\_squared
\end{minipage} & \begin{minipage}[b]{\linewidth}\raggedleft
p\_value\_slope
\end{minipage} & \begin{minipage}[b]{\linewidth}\raggedleft
F\_statistic
\end{minipage} \\
\midrule\noalign{}
\endhead
\bottomrule\noalign{}
\endlastfoot
value & Latitude & -1331.343 & 1.8718 & 0.0000 & 0.9751 & 0.001 \\
value1 & Longitude & -1333.361 & 15.8797 & 0.0034 & 0.2398 & 1.386 \\
\end{longtable}

\subsubsection{Scatter Plots with Regression
Statistics}\label{scatter-plots-with-regression-statistics}

\textbf{Purpose:} Creates scatter plots of latitude and longitude
against area difference, including regression lines and displaying R²
and p-values.

\subsection{Plot Accession Locations on a World
Map}\label{plot-accession-locations-on-a-world-map}

\textbf{Purpos}e: Prepares the world map data and creates a base plot
for plotting the locations of positive and negative accessions.

Plot the map and overlay with the accessions used in the GWAS

\begin{figure}[H]

{\centering \pandocbounded{\includegraphics[keepaspectratio]{01_Accession_phenotypes_files/figure-pdf/Fig1d-ACC-Geographic-1.pdf}}

}

\caption{Geographic distribution of accessions used in the GWAS.}

\end{figure}%

\section{Session information}\label{session-information}

\textbf{Purpose}: Print the session information, including details about
the R version, loaded packages, and the system environment. This is
useful for reproducibility, as it documents the computational
environment in which the analysis was performed.

\begin{verbatim}
R version: R version 4.4.1 (2024-06-14) 
\end{verbatim}

\begin{verbatim}
Attached packages:
\end{verbatim}

\begin{verbatim}
maps - 3.4.2.1 
dtw - 1.23-1 
proxy - 0.4-27 
fstcore - 0.9.18 
emmeans - 1.10.6 
lmerTest - 3.1-3 
lme4 - 1.1-35.5 
Matrix - 1.7-0 
broom - 1.0.7 
readxl - 1.4.3 
openxlsx - 4.2.8 
circlize - 0.4.16 
ComplexHeatmap - 2.22.0 
ggdendro - 0.2.0 
knitr - 1.49 
cowplot - 1.1.3 
MESS - 0.5.12 
ggpubr - 0.6.0 
DescTools - 0.99.58 
magrittr - 2.0.3 
gwaR - 0.6 
lubridate - 1.9.4 
forcats - 1.0.0 
stringr - 1.5.1 
dplyr - 1.1.4 
purrr - 1.0.2 
readr - 2.1.5 
tidyr - 1.3.1 
tibble - 3.2.1 
ggplot2 - 3.5.1 
tidyverse - 2.0.0 
\end{verbatim}




\end{document}
